\section{Panoramica di Sistema: Il Diagramma di Flusso delle 3 Fasi}
\subsection{Introduzione Architetturale} 
Il panorama dell'\textit{Automated Program Repair} è storicamente dominato da una dicotomia: 
da un lato, gli approcci basati su logica formale e analisi simbolica, eccellenti nella precisione 
diagnostica ma limitati nella capacità di sintetizzare codice complesso; dall'altro, i recenti sistemi 
basati su \textit{Large Language Models} (LLM), dotati di straordinarie capacità generative ma inclini 
a introdurre allucinazioni logiche e modifiche non semanticamente valide. In questo contesto, HybridRepair 
si configura come un sistema APR progettato attorno a un principio cardine dell'Ingegneria 
del Software: il rigoroso disaccoppiamento architetturale tra l'inferenza deterministica delle cause e la 
variabilità stocastica della generazione di codice.

Il termine "ibrido" definisce specificamente una natura neuro-simbolica: una componente simbolica fondata 
su LogicFL e Prolog viene delegata alla diagnosi strutturata e ripetibile del difetto, mentre una 
componente neurale (implementata tramite GPT-4 via Azure OpenAI) traduce tale diagnosi in conoscenza 
operativa e, infine, in codice sorgente corretto. Questa sinergia supera i limiti derivanti dai singoli 
approcci: la logica formale ancora il ragionamento probabilistico a verità inconfutabili del programma, 
riducendo drasticamente lo spazio di ricerca e garantendo che l'intervento del LLM sia sempre mirato e 
semanticamente fondato.

L'obiettivo principale del sistema è la riparazione in modalità \textit{fully-automated} di \textit{bug} NPE Java 
reali, estratti dal \textit{benchmark} standard di settore Defects4J. A partire da un identificatore di 
difetto (ad esempio, Codec-13), il sistema orchestra un'esecuzione attraverso tre fasi sequenziali, 
nettamente separate e disaccoppiate: FaultOracle, SpecReason e PatchAgent. Ognuna di queste fasi possiede 
responsabilità architetturali coese, interfacce di comunicazione tipizzate e confini ben delineati. Tale 
segregazione assicura che la tracciabilità del processo diagnostico rimanga sempre esplicabile 
(spiegabilità garantita dalle regole Prolog) e che il modello generativo riceva specifiche inequivocabili, 
mitigando alla radice l'ambiguità inerente al linguaggio naturale.


\subsection{\textit{Pipeline} dei dati: input e output delle tre fasi}
Per garantire che i moduli siano testabili in isolamento e sostituibili, l'intero flusso di informazioni
intra-fase è governato da un contratto di interfaccia esplicito. Le strutture dati che mediano le transizioni 
sono fortemente tipizzate e centralizzate nel modulo shared/models.py, realizzando un'architettura robusta e 
resistente agli errori di disallineamento dei dati.


La Tabella~\ref{tab:pipeline_io} riassume, per ciascuna delle tre fasi, le informazioni in ingresso, la
natura dell'elaborazione e l'artefatto prodotto. I paragrafi successivi ne descrivono il ruolo
concettuale; il funzionamento interno di ogni fase è trattato in dettaglio nelle sezioni dedicate
(\ref{sec:fase1}, \ref{sec:fase2} e \ref{sec:fase3}).

\begin{table}[H]
\centering
\small
\caption{Input, elaborazione e output delle tre fasi della \textit{pipeline}.}
\label{tab:pipeline_io}
\begin{tabularx}{\textwidth}{@{}l >{\raggedright\arraybackslash}X >{\raggedright\arraybackslash}X l@{}}
\toprule
\textbf{Fase} & \textbf{Input} & \textbf{Elaborazione} & \textbf{Output} \\
\midrule
\textbf{1. FaultOracle} & \texttt{bug\_id} e artefatti LogicFL pre-generati (\texttt{fault\_locs}, \texttt{root\_cause}, \texttt{stack\_traces}) & Analisi statica e deduzione logica (Prolog); nessun modello generativo & \texttt{FaultReport} \\
\addlinespace
\textbf{2. SpecReason} & \texttt{FaultReport} & Raccolta del contesto API + ragionamento \textit{Chain-of-Thought} con auto-critica; nessun codice prodotto & \texttt{RepairSpec} + \texttt{Ingredients} \\
\addlinespace
\textbf{3. PatchAgent} & \texttt{FaultReport}, \texttt{RepairSpec}, \texttt{Ingredients} & Agente LLM iterativo (FSM) che genera, applica, compila e testa la \textit{patch} in \textit{sandbox} & \texttt{RepairResult} \\
\bottomrule
\end{tabularx}
\end{table}

\paragraph{Fase 1 --- FaultOracle.} Ha l'onere esclusivo di localizzare il difetto nel codice sorgente
Java e di derivarne una diagnosi causale strutturata. L'operazione esclude categoricamente i modelli
generativi, affidandosi solo ad analisi statica e ragionamento deduttivo: l'analisi causale
deterministica deve sempre precedere quella generativa, per evitare che l'LLM ``inventi'' problemi
inesistenti. L'esito è un oggetto \texttt{FaultReport} che aggrega le locazioni del \textit{fault}
ordinate per sospiciosità, le catene causali (grezze e semanticamente ancorate) e il contesto dinamico
del fallimento.

\paragraph{Fase 2 --- SpecReason.} Opera la transizione dal dominio formale della diagnosi a quello
semantico-generativo: la diagnosi di Fase~1 viene trasformata in una specifica di riparazione in
linguaggio naturale attraverso un ragionamento \textit{Chain-of-Thought}. In questa fase
\textit{non si produce ancora codice}; l'obiettivo è la sintesi di una comprensione ragionata del problema,
ancorata alle sole API realmente disponibili e validata da un passo di auto-critica. L'output è la coppia
\texttt{RepairSpec} + \texttt{Ingredients}.

\paragraph{Fase 3 --- PatchAgent.} Materializza la comprensione in codice eseguibile. Un agente LLM,
governato da una Macchina a Stati Finiti, genera, applica, compila e testa la \textit{patch} operando
su una \textit{sandbox} isolata del sorgente, con un meccanismo di \textit{retry} adattivo che apprende
dai fallimenti precedenti. Il processo culmina nell'oggetto \texttt{RepairResult}, che documenta l'esito,
il tentativo vincente e lo storico completo dei tentativi.